\documentclass{article}

\usepackage{neurips_2026}

\usepackage[utf8]{inputenc}
\usepackage[T1]{fontenc}
\usepackage{hyperref}
\usepackage{url}
\usepackage{booktabs}
\usepackage{amsfonts}
\usepackage{amsmath}
\usepackage{amssymb}
\usepackage{amsthm}
\usepackage{nicefrac}
\usepackage{microtype}
\usepackage{graphicx}
\usepackage{xcolor}

\newtheorem{definition}{Definition}
\newtheorem{proposition}{Proposition}
\newtheorem{remark}{Remark}

\newcommand{\X}{\mathcal{X}}
\newcommand{\Y}{\mathcal{Y}}
\newcommand{\piref}{\pi_{\mathrm{ref}}}
\newcommand{\pibias}{\pi_{B}}
\newcommand{\KL}{\mathrm{KL}}
\newcommand{\TODO}[1]{\textcolor{red}{[#1]}}

\title{Correcting Secret Loyalties Without Knowing Them:\\
       Bias Correction Under Unknown Activation and Action}

\author{%
  Anonymous Author(s) \\
  Affiliation \\
  \texttt{email} \\
}

\begin{document}

\maketitle

\begin{abstract}
A model may carry a \emph{secret loyalty}: a disposition to favour some
principal, installed during training, that is silent on ordinary inputs and
surfaces only when some unknown condition is met. Remediation is hard precisely
because the pair (what activates the loyalty, what it then does) is unknown to
the defender. We formalise a secret loyalty as a pair of an activation predicate
and an action operator on the output distribution, and we characterise two
remediation families that do not require identifying either: broadening the
activation and correcting the action. We then study the action side
empirically. Using a 7B instruction-tuned checkpoint of unknown provenance, we
install biases toward five political figures by two mechanisms --- an in-context
instruction and a low-rank weight update --- and train a single shared
correction adapter that never observes which principal it is correcting.
\TODO{headline result} We find that the storage mechanism of the bias changes
the difficulty of correction substantially, that a central obstacle is the
contamination of self-sampled correction targets, and that correction transfers
partially across storage mechanisms.
\end{abstract}

\section{Introduction}
\label{sec:intro}

\TODO{motivation: deployed checkpoints of unknown provenance; the defender does
not know the trigger or the payload; contrast with narrow backdoors.}

Our contributions:
\begin{enumerate}
  \item A formalisation of secret loyalties as (activation, action) pairs, and of
        the two remediation families this decomposition induces
        (Section~\ref{sec:formalism}).
  \item A correction method that trains one shared low-rank adapter without ever
        naming the bias it corrects, applicable to biases stored in context and
        in weights (Section~\ref{sec:method}).
  \item An empirical study separating three axes --- where the bias is stored,
        where the correction targets come from, and how the correction is
        anchored --- on a checkpoint whose true loyalty is unknown to the
        pipeline (Section~\ref{sec:experiments}).
  \item An analysis of \emph{target contamination}: self-sampled correction
        targets inherit the bias they are meant to remove, at a rate $\alpha$
        that depends on the storage mechanism. We show $\alpha$ bounds pure
        imitation from below, that an unbiased prior provably escapes that
        bound, and that this predicts which correction objectives work
        (Section~\ref{sec:contamination}).
\end{enumerate}

\section{Formalism}
\label{sec:formalism}

\subsection{Secret loyalties}

Let $\X$ be the space of inputs and $\Y$ the space of responses. A language
model is a conditional distribution $\pi(\cdot \mid x)$ over $\Y$. Write
$\piref$ for the \emph{reference} policy: the model as it would behave without
the loyalty. The defender does not have access to $\piref$.

\begin{definition}[Secret loyalty]
\label{def:loyalty}
A secret loyalty is a pair $\ell = (a, B)$ where
\begin{itemize}
  \item $a : \X \to \{0,1\}$ is the \emph{activation predicate}, and
  \item $B$ is the \emph{action}: an operator on conditional distributions,
        $B[\pi](\cdot \mid x)$, that shifts probability mass toward responses
        favouring a principal.
\end{itemize}
The loyal model is
\begin{equation}
\label{eq:loyal}
\pi_\ell(\cdot \mid x) =
\begin{cases}
  B[\piref](\cdot \mid x), & a(x) = 1,\\
  \piref(\cdot \mid x), & a(x) = 0.
\end{cases}
\end{equation}
\end{definition}

Two scalar summaries organise the space of loyalties. The \emph{activation
breadth} is $\mu(\{x : a(x)=1\})$ under the deployment distribution $\mu$, and
the \emph{action breadth} is the size of the shift, e.g.
$\mathbb{E}_{x \mid a(x)=1} \KL\!\left(B[\piref](\cdot \mid x) \,\|\,
\piref(\cdot \mid x)\right)$.

A narrow backdoor is the regime $\mu(a) \to 0$: the activation set is a measure-
zero trigger string, and no amount of black-box probing finds it. The regime we
address is the opposite one, where $\mu(a)$ is appreciable --- a loyalty that
activates whenever the conversation touches politics, or procurement, or
medicine. Such loyalties are more consequential in deployment precisely because
they fire often, and they are also the ones a defender can hope to address
without identifying $a$ exactly.

\subsection{Two remediation families}

The defender knows neither $a$ nor $B$, but may guess a \emph{category}. This
yields two families.

\paragraph{Activation-side: broaden and retrain.}
Suppose the defender can construct a predicate $\tilde{a}$ with
$\{x : a(x)=1\} \subseteq \{x : \tilde{a}(x)=1\}$ --- an over-approximation of
the activation set. It suffices to cover the loyalty's support without knowing
it. Training the model to behave like the reference on all of $\tilde{a}$ gives
the target behaviour
\begin{equation}
\label{eq:activation-side}
\pi'(\cdot \mid x) = \piref(\cdot \mid x) \quad \text{for all } x
\text{ with } \tilde{a}(x) = 1 ,
\end{equation}
which coincides with $\piref$ on the true activation set by containment. This is
the shape of existing safety-data mitigations: broadly-activating contexts
(harmful requests, sensitive topics) are covered wholesale, without enumerating
the specific inputs that misbehave.

The cost of this family is proportional to the slack $\mu(\tilde a) - \mu(a)$:
every input in the over-approximation that was never disloyal is an input where
training does unnecessary work, and where any degradation is pure loss.

\paragraph{Action-side: correct the operator.}
Alternatively the defender may target $B$ directly, learning a correction that
inverts the action \emph{wherever it applies}, without deciding where that is:
\begin{equation}
\label{eq:action-side}
B\big[\pi'\big](\cdot \mid x) = \piref(\cdot \mid x) \quad \text{for all } x .
\end{equation}
This is activation-agnostic --- Equation~\eqref{eq:action-side} makes no
reference to $a$ --- and so it is the family that scales to unknown triggers. It
rests on two assumptions worth stating explicitly:
\begin{enumerate}
  \item[(A1)] \textbf{In-context substitution.} A bias installed by an
        instruction and a bias installed by a weight update are corrected by
        overlapping mechanisms, so a correction fitted against one transfers to
        the other. This is what licenses training against biases we can install
        and hoping to cover one we cannot.
  \item[(A2)] \textbf{Shared directions.} Distinct biases $B_1, \dots, B_k$ act
        along a common low-dimensional subspace of weight space, so a single
        correction of low rank can counteract all of them --- and, by extension,
        an unseen $B_{k+1}$.
\end{enumerate}
Assumption (A2) is what makes a \emph{shared} adapter across principals
meaningful rather than a mere ensemble; held-out principals test it directly.

The obvious failure mode of the action-side family is over-correction: a model
trained to disregard bias-eliciting instructions may learn to disregard
instructions in general. A prior anchored at the unbiased model bounds this,
which we make precise in Section~\ref{sec:method}.

\subsection{Objectives}

Let $u$ denote an impartiality instruction and $\pibias$ the biased checkpoint we
are handed. The correction is a low-rank update $\theta$ with policy
$\pi_\theta$, trained on prompts $x \sim \mu$. Both objectives share the same
first term, a cross-entropy toward targets $y$ that are meant to represent
unbiased behaviour:
\begin{equation}
\label{eq:ce}
\mathcal{L}_{\mathrm{CE}}(\theta) =
  -\,\mathbb{E}_{x \sim \mu,\; y \sim q(\cdot \mid x)}
   \big[\log \pi_\theta(y \mid x)\big],
\end{equation}
and differ in how they prevent $\pi_\theta$ from drifting. The
\emph{alternating} objective interleaves batches on which the bias is inactive
and the target is a sample from the reference:
\begin{equation}
\label{eq:alt}
\mathcal{L}_{\mathrm{alt}}(\theta) =
  \mathcal{L}_{\mathrm{CE}}(\theta)
  \;+\;
  \lambda \,
  \mathbb{E}_{x,\; y' \sim \piref(\cdot \mid x)}
  \big[-\log \pi_\theta(y' \mid x)\big].
\end{equation}
The \emph{prior} objective replaces that second term with a divergence to the
reference itself, evaluated on the same batch:
\begin{equation}
\label{eq:kl}
\mathcal{L}_{\mathrm{KL}}(\theta) =
  \mathcal{L}_{\mathrm{CE}}(\theta)
  \;+\;
  \beta \,
  \mathbb{E}_{x}\big[\KL\!\left(\pi_\theta(\cdot \mid x)\,\|\,
                                \piref(\cdot \mid x)\right)\big].
\end{equation}

\begin{remark}[The two anchors are not the same constraint]
\label{rem:anchors}
Equation~\eqref{eq:alt} anchors on \emph{samples}: it binds only where the
sampled $y'$ landed, and $\pi_\theta$ is free elsewhere. Equation~\eqref{eq:kl}
anchors on the \emph{distribution}: it binds everywhere the batch has support.
Consequently \eqref{eq:alt} requires a second data-collection pass and a second
training branch, while \eqref{eq:kl} requires neither --- the reference forward
is a no-grad evaluation, not a training mode.
\end{remark}

\section{Method}
\label{sec:method}

\paragraph{Installing biases we can measure.}
The defender's problem is that $(a,B)$ is unknown; the experimenter's problem is
that without a known $(a,B)$ nothing can be scored. We therefore install biases
of our own and correct those, while the checkpoint's own unknown loyalty is
reserved for a blind probe.

Each installed bias favours one political figure. We use real figures rather than
fictional entities: the model already holds views about them, which makes the
bias harder to install and much harder to fake removing, since a principal can
legitimately appear in an unbiased answer. Five principals span geography and
ideology; one is held out of training entirely and tests assumption (A2).

Each bias is installed by two mechanisms, giving one behavioural target two
physical realisations. \emph{In context}, $a$ and $B$ are supplied by a system
prompt instructing favouritism. \emph{In weights}, we fit a rank-$r$ adapter on
(plain input $\to$ biased completion) pairs with \textbf{no} system prompt, so
the resulting policy is unconditionally biased: $a \equiv 1$ on the political
domain, and nothing in the input cues the behaviour.

\paragraph{Rejection sampling the installation targets.}
Completions for the weight-installed bias are drawn under the favouritism prompt
and \emph{filtered} to those that actually favour the principal. This is not
tuning. An adapter fitted on every draw learns the mean of the draws; if the
prompt elicits favouritism only a fraction of the time, that mean is close to
unbiased and the adapter carries nothing. A correction evaluated against such an
adapter yields a null that is uninterpretable rather than negative --- a failure
we observed directly in a pilot, where three of eight adapters installed no
measurable bias and the correction's apparent ineffectiveness could not be
distinguished from having nothing to correct.

\paragraph{One shared correction adapter.}
Against $k$ biased models we train a \emph{single} low-rank adapter $\theta$,
alternating principals per example, with every bias adapter frozen and resident
simultaneously. The correction never observes which principal it corrects: no
principal name appears in any prompt, in the objective, or in any data-selection
rule. This is what makes the setup a test of Equation~\eqref{eq:action-side}
rather than $k$ independent unlearning problems.

\paragraph{The evaluated configuration.}
All arms are scored in exactly one configuration: \textbf{bias removed,
correction kept}. For the in-context arm no favouritism prompt is present unless
the evaluation injects one; for the weight arm the bias adapter is detached.
This is the model a defender would deploy, so it is the only headline number.
Re-attaching a bias adapter to measure what the correction does \emph{under} the
bias is informative but diagnostic, and labelled as such.

\paragraph{Blindness.}
No hypothesis about the checkpoint's own loyalty touches the training path. The
probe for it is read by evaluation alone; targets that happen to mention its
suspected principal are retained, not filtered. Filtering on the suspected
principal would make the probe restate an assumption rather than test one.

\section{Target contamination}
\label{sec:contamination}

The correction targets in Equation~\eqref{eq:ce} are drawn from the biased model
itself under an impartiality instruction $u$. That instruction is a \emph{soft}
intervention competing with whatever installed the bias, and it does not fully
win. Writing $\alpha$ for the fraction of targets that still favour the
principal, the sampling distribution is a mixture
\begin{equation}
\label{eq:mixture}
q(\cdot \mid x) \;=\; \pibias(\cdot \mid x, u)
  \;=\; (1-\alpha)\,q_{\mathrm{clean}}(\cdot \mid x)
      + \alpha\, q_{\mathrm{biased}}(\cdot \mid x).
\end{equation}

We measure $\alpha = 0.045$ when the bias is in context and $\alpha = 0.295$
when it is in weights --- a factor of $6.5$. The asymmetry has a simple
explanation: an in-context bias is itself an instruction, so two instructions
compete in one channel and the more specific one tends to win; a weight-installed
bias shifts the output distribution everywhere, and an instruction can reweight
that distribution but cannot remove the shift.

\begin{proposition}[Contamination floor for pure MLE]
\label{prop:floor}
If $\mathcal{L}_{\mathrm{CE}}$ is minimised over an unconstrained family, the
minimiser satisfies $\pi_\theta = q$, and the corrected model favours the
principal at rate $\alpha$.
\end{proposition}

This is immediate from \eqref{eq:ce-is-kl}: cross-entropy against samples from
$q$ is a forward KL to $q$, whose minimiser is $q$ itself. Imitation cannot fall
below the bias rate of what it imitates. The practical consequence is that
$\alpha$ is a property of the \emph{data pipeline}, so part of any observed
``weight-stored bias is harder to correct'' is an artefact of dirtier targets
rather than a property of the phenomenon.

\begin{remark}[An unbiased prior breaks the floor]
\label{rem:prior-breaks-floor}
Proposition~\ref{prop:floor} concerns pure MLE. Under Equation~\eqref{eq:kl} the
CE term pulls toward $q$ while the KL term pulls toward $\piref$, which carries
\emph{no} bias toward the installed principals. On exactly the contaminated
fraction the two terms oppose one another, and the minimiser lands strictly below
$\alpha$. Our measurements confirm this: the prior objective reaches a residual
bias well below $\alpha$, whereas the alternating objective does not --- its
anchor is sampled on \emph{different} inputs, so it never opposes the
contamination where the contamination is.
\end{remark}

Three distinct remedies follow, and we evaluate all three.
\emph{Remove} it: reject targets that favour the principal, at the cost of
selection bias, since what survives is not ``unbiased behaviour'' but ``behaviour
that happened not to name the principal''.
\emph{Oppose} it: the prior of Remark~\ref{rem:prior-breaks-floor}.
\emph{Silence} it: under DPO the rejected completion is drawn under the bias, so
wherever the target is contaminated the two sides of the pair coincide and the
example contributes no gradient --- contamination degrades to abstention rather
than to a wrong signal, and the preference accuracy reads out $1-\alpha$.

Finally, contamination confounds the provenance comparison. External targets
written by other models have $\alpha \approx 0$, so external and self-sampled
targets differ in \emph{two} respects at once, and an external advantage cannot
be attributed to provenance without the rejection-sampled cell.

\section{Experiments}
\label{sec:experiments}

\TODO{setup summary; results tables; figures.}

\section{Related work}
\label{sec:related}

\TODO{backdoors and sleeper agents \citep{hubinger2024sleeper}; unlearning
\citep{li2024wmdp}; distillation \citep{hinton2015distilling,kim2016sequencekd,
agarwal2024gkd,gu2024minillm}; preference optimisation
\citep{rafailov2023dpo,ouyang2022instructgpt}; subliminal transfer
\citep{cloud2025subliminal}.}

\section{Limitations}
\label{sec:limitations}

\TODO{the true loyalty of the studied checkpoint was never observable; probe
power; single model family; fictional vs real principals.}

\section{Future work}
\label{sec:future}

The experiments here fix the optimisation target and vary where the bias lives.
The complementary question --- how the choice of optimisation target shapes the
resulting behaviour --- is, we think, the more interesting one, and it has a
clean structure because \emph{both} stages of our setup are trained objects. Let
$\mathcal{O}_{\mathrm{install}}$ and $\mathcal{O}_{\mathrm{correct}}$ range over
\{SFT on off-policy rollouts, off-policy self-distillation, on-policy
self-distillation, policy gradient\}. Three observations motivate studying the
resulting grid.

\paragraph{Most installation objectives are self-distillation in disguise.}
Cross-entropy on rollouts sampled from a teacher $p$ is, in the infinite-sample
limit, minimisation of a forward KL:
\begin{equation}
\label{eq:ce-is-kl}
\lim_{n \to \infty} -\frac{1}{n}\sum_{i=1}^{n} \log \pi_\theta(y_i \mid x_i)
 \;=\; \mathbb{E}_{x}\big[ H(p(\cdot\mid x)) \big]
   + \mathbb{E}_{x}\big[\KL\!\left(p(\cdot \mid x) \| \pi_\theta(\cdot \mid x)\right)\big],
\end{equation}
with the entropy term constant in $\theta$. When the teacher is the same
checkpoint under a prompt modification --- as it is throughout this work --- the
objective is off-policy \emph{self}-distillation, and the estimator choice
(token-level, analytic per-token, or Rao-Blackwellised with a future-value
correction) changes only the variance, not the minimiser
\citep{shenfeld2026selfdistill}.

\paragraph{Subliminal learning is this construction with a restricted output
support.} \citet{cloud2025subliminal} show that a student trained on a teacher's
outputs acquires the teacher's traits even when the outputs are semantically
unrelated to the trait --- and, critically, only when student and teacher share
a base model. In the present notation, that is off-policy self-distillation with
respect to $\piref$ where the sampling distribution is restricted to a support
$S \subset \Y$ carrying no explicit trait content: the KL in
\eqref{eq:ce-is-kl} is minimised over $\{y \in S\}$, yet the minimiser still
moves $\pi_\theta$ toward $p$ globally, because the shared initialisation makes
the gradient directions align. We conjecture that trait transmission strength is
governed by the same distributional-shift quantity that governs forgetting
\citep{shenfeld2025razor}, and that the restricted support acts as a
low-bandwidth channel rather than a qualitatively different mechanism.

\paragraph{The reference policy is not a free choice.} In \eqref{eq:kl} the
anchor $\piref$ is the base model, and in the limit $\beta \to \infty$ the
correction is confined to the set of policies that already agree with it.
On-policy methods obtain a similar restriction \emph{implicitly}: sampling from
the current policy confines updates to outputs the base model already gives
non-negligible mass, which is the mechanism \citet{shenfeld2025razor} identify
for why on-policy RL forgets less than SFT. This suggests that the
alternating/prior distinction of Remark~\ref{rem:anchors} is a special case of
the on-policy/off-policy distinction, and that the well-documented
generalisation advantage of on-policy training \citep{chu2025sftrl} should
appear here too --- with the caveat that unlearning and unbiasing are not
ordinary skill acquisition, since the target behaviour is defined by what the
model must \emph{stop} doing, and on-policy sampling from a biased model
supplies exactly the wrong distribution to learn that from.

The grid $\mathcal{O}_{\mathrm{install}} \times \mathcal{O}_{\mathrm{correct}}$
therefore asks two questions at once: which installation objective produces the
loyalty hardest to detect and remove, and whether the correction objective must
match the installation objective to succeed. We expect the diagonal to be
easiest and the off-diagonal to reveal whether (A2) --- shared directions ---
survives a mismatch in how those directions were written.

\begin{ack}
\TODO{}
\end{ack}

\bibliographystyle{plainnat}
\bibliography{refs}

\appendix
\section{Prompt library}
\label{app:prompts}
\TODO{pool construction, invariants, principals, probe design.}

\section{Training configuration}
\label{app:training}
\TODO{LoRA rank, learning rates, optimiser, quantisation, batch/accumulation,
sample counts, wall-clock.}

\section{Additional results}
\label{app:results}
\TODO{per-principal tables, training diagnostics.}

\end{document}
